\abstract
\justify
In aerial robotics, shared control is widely exploited in contexts where high-level decision-making capabilities provided by the human are required and, at the same time, robot can assist human operator in performing the task in a safe and efficient way. In some cases, like complex inspections, the execution of semi-autonomous tasks via aerial robotics is obstructed by the limited situational awareness provided by onboard first-person view cameras. For this reason it could be useful to use two UAVs (Unmanned Aerial Vehicles): one (Inspection Drone) effectively acting for the inspection task, the other one (Observer Drone) as visual aid for the human operator. In this context, teleoperation can been used to let the operator intervene on both drones for actively monitoring the navigation and inspection, so to improve general safety and task performance. Because of all difficulties that arise from teleoperating two drones, the architecture features a shared-control paradigm. Precisely, this thesis proposes a human-robot shared control framework utilizing a two-UAVs architecture, exploiting \textit{exocentric vision} (i.e., third person view) to feed the human with the "best" view of the scene.
\\To achieve this, the observation task of Observer UAV is formulated as a nonlinear programming problem and solved via Nonlinear Model Predictive Control (NMPC). The NMPC controller minimizes the geometric 3D tracking error to keep observed UAV in the field of view at a desired relative distance, with a desired point Of view. Moreover, to prevent operator motion sickness, the controller explicitly minimizes acceleration, jerk and snap, ensuring smooth camera motion while respecting actuation constraints (thrust and torques). The framework features bilateral teleoperation: if human commands on Observer UAV try to crash it over the other drone or loose it from the view, a force feedback based on repulsive artificial potential field method is generated to guide the operator, without mathematically altering human command. Simultaneously, the Inspection UAV can perform an inspection task executing the contact utilizing an admittance filter planner. Human operator can perform the interaction through teleoperation, relying on the admittance filter that provides compliance to the drone during contact. 
\\Simulations to validate this architecture are conducted in Gazebo, using ROS2 middleware and PX4 firmware on the two drones. Simulation results demonstrate the real-time computational feasibility of the NMPC solver at 100 Hz and highlight a bounded geometric tracking error, ensuring the target is never lost from the FOV even under conflicting human inputs. The framework provides contact stability for the Inspection UAV and proves the effectiveness of haptic feedback in safely integrating human high-level decision-making with autonomous low-level constraints. 
\\Preliminary experiments have been conducted, validating simulation results only  over the Observer UAV, showing the effectiveness of NMPC controller and bilateral teleoperation in real scenarios, with a light sim-to-real gap. 